\section{The objects of the coset system}
Use the three rooted cores A, B and C and the notation of the attachment
theorem. Thus $Q=\{1,2,3,4\}$, $H=\langle h_0,h_1\rangle$, and a tree
$T$ meets $Q$ at its leaf $d$. Its unique edge at $d$ is denoted $b$.
The group of the complete marking is $G=\Sym(\Omega)$.
For each generator $i$, let $G_i$ be the subgroup generated by all the
other generators. Elements of type $i$ are left cosets of $G_i$;
incidence means nonempty intersection. We work with images of the
following objects under permutations of $\Omega$.

For an edge $e$ of $T$, let $S_e$ be the component of $T-e$ not containing
$d$. For $e\ne b$, the attachment formula gives
\[
G_e=\Sym(S_e)\times\Sym(\Omega\setminus S_e).
\]
Thus an object of type $e$ is a subset $X_e\subset\Omega$ of size $|S_e|$.
The sets $S_e$ form a laminar family: two of them are disjoint or one
contains the other. Their images have the same pairwise incidence rule.
Indeed, a permutation carries two prescribed subsets to two given subsets
if and only if their four Boolean regions have the prescribed sizes.
In the nested or disjoint cases this is precisely the stated containment
or disjointness relation, together with the fixed individual sizes.

The bridge stabilizer is $G_b=H\times\Sym(\Omega\setminus Q)$.
A bridge object is a four-set $Q'$ together with the following decoration.
In A and B it is a partition $M$ of $Q'$ into two unordered pairs,
the image of $\{\{1,4\},\{2,3\}\}$. Its stabilizer on $Q'$ is $D_8=H$.
In C it is an ordered list of the three perfect matchings on $Q'$,
the image of a fixed order of those on $Q$.
The action of $S_4$ on these three matchings is onto $S_3$ and has
kernel $\{1,(12)(34),(13)(24),(14)(23)\}=H$.
This describes the bridge cosets exactly, including their full multiplicity.

The two remaining objects, of types 0 and 1, are as follows.
Every displayed tuple has ordered entries; a displayed pair such as
$\{a,b\}$ is unordered.
\[
\begin{array}{c|cc}
 &\text{type }0&\text{type }1\\\hline
 A& P=\{1,2\}&(p,R)=(1,\{2,3\})\\
 B& P=\{1,2\}&(p,q)=(1,4)\\
 C& P=\{1,3\}&R=\{1,2\}
\end{array}
\]
These are exact coset coordinates, because the attachment formula gives
their stabilizers as $G_0$ and $G_1$. In A, $p\notin R$; in B,
$p\ne q$; in C each object is a two-subset. All their displayed points
will be called their support.

Two such core objects are incident precisely under the following rules:
\begin{equation}\label{eq:coreflags}
\begin{array}{c|l}
 A&p\in P,\quad |P\cap R|=1,\quad p\notin R;\\
 B&p\in P,\quad q\notin P;\\
 C&|P\cap R|=1.
\end{array}                                                    \end{equation}
For example, in A these conditions identify the images of 1, 2 and 3
as $p$, the unique point of $P\cap R$, and the other point of $R$.
The fourth core point may be any unused point. The other two assertions
follow in the same way. In each case a compatible pair has support of
size three; a single core object has support of size at most three.

A core object is incident with an ordinary tree object $X_e$ precisely
when its support is disjoint from $X_e$. Its stabilizer has one large
block, containing the whole tail; that block is the complement of its
support. Necessity follows by taking images of the standard objects.
Conversely, disjointness lets one map the displayed core points first,
then map $S_e$ to $X_e$, and extend these disjoint partial bijections.
The prescribed cardinalities guarantee enough remaining points.
Similarly a bridge object $(Q',\delta)$ is incident with $X_e$ precisely
when $Q'\cap X_e=\varnothing$. These arguments apply also when the
support has several distinguished blocks: the internal labels are fixed
by the core-object choice and place no condition on the tail.

\section{The finite decoration check}
\begin{lemma}\label{lem:decoration}
A flag consisting of a bridge object and any subset of the two core
types has a representative bijection $Q\longrightarrow Q'$.
\end{lemma}
\begin{proof}
Incidence with the bridge puts each support inside $Q'$ and imposes
its orbit under $H$. It remains to check simultaneous realization when
both core types occur.

For A, a type-0 pair $P$ is a transversal of the matching $M$.
A type-1 object consists of one pair $R$ of $M$ and a point $p$ of the
other pair. By \eqref{eq:coreflags}, $p$ is the point of $P$ outside $R$, and $P\cap R$
has one point. Map 1 to $p$, 2 to that intersection point, 3 to the
other point of $R$, and 4 to the remaining point. This maps the
standard matching and both standard core objects as required.

For B, incidence with the matching says that $P$ is a transversal and
$(p,q)$ is an ordering of one of its pairs. Rule \eqref{eq:coreflags} puts $p$ in $P$
and $q$ outside $P$. Map 1 to $p$, 4 to $q$, 2 to the other point
of $P$, and 3 to the remaining point. Again the three objects are
simultaneously realized.

For C, $P$ is one pair of the matching assigned to type 0 and $R$ is
one pair of the different matching assigned to type 1. Distinct
perfect matchings on a four-set have every pair meeting every pair
in one point. Map 1 to $P\cap R$, 2 to the other point of $R$,
3 to the other point of $P$, and 4 to the remaining point. The first
two ordered matchings are correct; the third is then determined.
This is a representative of the given bridge decoration as well.

If only one core type occurs, any incident representative already
provides the required bijection. The empty case is immediate.
\end{proof}

\section{Completion of a tree flag}\label{sec:regularity}
We give explicitly the weighted-tree fact used here. Regard the root
$d$ as a vertex of weight four, occupied by the whole set $Q$, and every
other tree vertex as having weight one. For selected edges $e\ne b$,
suppose subsets $X_e$ of the prescribed sizes have all their prescribed
pairwise nesting and disjointness relations. For a selected edge $e$,
let its selected children mean the maximal selected descendants below it.
The differences
\[
 X_e\setminus\bigcup_{f\text{ a selected child of }e}X_f
\]
have exactly the cardinalities of the analogous differences among the
standard $S_e$. The root region
\[
 D=\Omega\setminus\bigcup_{e\text{ maximal selected}}X_e
\]
has the cardinality of the analogous standard root region, and hence
has at least four points. All these regions are disjoint. A bijection
on a chosen four-subset of $D$ can therefore be extended independently
on each difference region to a permutation taking every $S_e$ to $X_e$.
This is the classical laminar flag-completion argument, with root
weight four, written in full; no classification theorem is needed for it.

\begin{theorem}\label{thm:flags}
For each of the three cores, and every attached tree allowed by the
attachment theorem, the coset system of the complete marking is
flag-transitive. It is a thin, residually connected geometry, and the
specified symmetric group acts regularly on its chambers.
\end{theorem}
\begin{proof}
Take an arbitrary flag and collect its ordinary tree objects $X_e$.
They have the nesting and disjointness relations above, so construct
the root region $D$. Every core-object support is contained in $D$,
because it is disjoint from every selected $X_e$.

If the flag contains a bridge object, its four-set $Q'$ also lies in
$D$. Lemma~\ref{lem:decoration} gives a bijection $Q\to Q'$ realizing
all the flag's core and bridge objects. If the bridge is absent, then
there are at most two core objects. Pairwise incidence gives \eqref{eq:coreflags}, and
hence a partial bijection from at most three core points to their
support. Extend it to a bijection $Q\to Q'$ for some four-subset
$Q'$ of $D$. This is possible because $|D|\geq4$. With just one or no
core object the same construction is simpler.

Now extend that bijection by the laminar completion just proved. The
resulting permutation realizes every object in the flag simultaneously.
It belongs to all of the corresponding cosets, so every flag extends
to a chamber and every flag of fixed type is in the orbit of the
standard flag. Thus the coset system is flag-transitive.

For completeness, the attachment theorem identifies every intersection
of maximal parabolics with the subgroup generated by the remaining
marks. The intersection of all maximal parabolics is trivial, so the
chamber action is free. A panel stabilizer is the cyclic group generated
by its omitted involution and has order two; hence panels have exactly
two chambers and the geometry is thin. In any residue of rank at least
two, the remaining maximal parabolics generate its group, since each
remaining generator belongs to at least one of those parabolics.
Indeed, for a coset system, the union of the subgroups specifying
its types generates a group precisely when its incidence graph is
connected: a word in those subgroups gives a path through successive
cosets, and conversely such a path gives a word. Applying this to each
residue proves residual connectedness. The chamber graph is connected
as well, since the involutions generate $G$. A type-preserving
automorphism fixing one chamber must fix its unique adjacent chamber
of each type, by thinness, and hence every chamber. It fixes every
object, so the full type-preserving automorphism group is exactly $G$.
\end{proof}

